\begin{pastexam}[label={px:ncku-106-cce-ds-08}]{成大電通所 106 學年度 第三大題第 1 題}
（填充；10 分）

Give the following algorithm whose input is a graph G, an edge weight w, and a vertex r in G.
\begin{quote}
WeightGraph(G, w, r)\\
\hspace*{1em}for each u $\in$ V(G)\\
\hspace*{2em}key[u] = $\infty$\\
\hspace*{1em}Insert all nodes into Q; \quad // Q denotes a priority queue.\\
\hspace*{1em}key[r] = 0;\\
\hspace*{1em}while (Q is not empty)\\
\hspace*{2em}u = ExtractMin(Q)\\
\hspace*{2em}for each neighbor v of u and v $\in$ Q\\
\hspace*{3em}if(\underline{\hspace{2em}A\hspace{2em}})\\
\hspace*{4em}\underline{\hspace{3em}B\hspace{3em}}
\end{quote}
\begin{enumerate}[label=\alph*.]
\item \textbf{(4 pts)} Please list the answers of A and B if the algorithm is to find the minimum
spanning tree.
\item \textbf{(4 pts)} Please list the answers of A and B if the algorithm is to find the single
source shortest tree.
\item \textbf{(1 pts)} Please give the time complexity in term of V and E.
\item \textbf{(1 pts)} Please give the time complexity if we replace Q by an array A.
\end{enumerate}
\end{pastexam}

\begin{pastsoln}
這個框架就是 Prim 和 Dijkstra 共同的樣子：每次從 $Q$ 取出 key 最小的頂點 $u$，再用 $u$ 的邊去更新
鄰居的 key。 兩者只差在 key 的意義。 要得到樹，B 除了改 key，也要記下 $v$ 的父節點
$\pi[v] = u$（框架裡沒有宣告 $\pi$，這裡補上；只填 key 的更新也算答對主要部分）。

\textbf{a.} 最小生成樹（Prim，資結書第~24.3~節）：$\mathit{key}[v]$ 是「把 $v$ 接到目前這棵樹上
最便宜的邊」的權重。
\begin{itemize}
\item A：\texttt{w(u, v) < key[v]}
\item B：\texttt{key[v] = w(u, v); }$\pi[v] = u$（優先佇列做 DecreaseKey）
\end{itemize}

\textbf{b.} 單一起點最短路徑樹（Dijkstra，資結書第~25.2~節）：$\mathit{key}[v]$ 是目前找到從 $r$ 到
$v$ 的最短距離。 這就是鬆弛（資結書虛擬碼~25.2）。
\begin{itemize}
\item A：\texttt{key[u] + w(u, v) < key[v]}
\item B：\texttt{key[v] = key[u] + w(u, v); }$\pi[v] = u$
\end{itemize}
Dijkstra 需要邊的權重非負。 框架只看還在 $Q$ 裡的鄰居，對 Dijkstra 也沒有問題：已經取出的頂點
距離已經是最短的，不會再被改小（資結書定理~25.6）。

\textbf{c.} $Q$ 用二元堆積：ExtractMin $V$ 次、每次 $O(\log V)$；DecreaseKey 最多 $E$ 次（無向圖每條
邊從兩端各看一次，$2E$ 次）、每次 $O(\log V)$；建堆積 $O(V)$。 總共 $O((V + E) \log V)$，連通圖時
就是 $O(E \log V)$。 （用 Fibonacci 堆積可以降到 $O(E + V \log V)$。）

\textbf{d.} 改用陣列存 key：ExtractMin 要掃整個陣列找最小值，$O(V)$，做 $V$ 次共 $O(V^2)$；更新
key 是直接改陣列的一格，$O(1)$，共 $O(E)$。 總共 $O(V^2 + E) = O(V^2)$。 稠密圖（$E$ 接近 $V^2$）
時這反而比堆積版的 $O(V^2 \log V)$ 快。
\end{pastsoln}
